% ---------------------------------------------------------------------------
%  FTEC5660 Reproducibility Work -- Mem0 on LOCOMO
%  Xiuling PIAO (1155267871)
%
%  Compile with pdfLaTeX, twice (the footer needs \pageref{LastPage}):
%      pdflatex REPORT.tex
%      pdflatex REPORT.tex
%  On Overleaf: upload REPORT.tex, set the compiler to pdfLaTeX.
%
%  If the PDF runs past the 10-page limit, change the first line to 10pt
%  and/or narrow the margins in \geometry.
% ---------------------------------------------------------------------------
\documentclass[11pt,a4paper]{article}

\usepackage[T1]{fontenc}
\usepackage[utf8]{inputenc}
\usepackage{newtxtext,newtxmath}
\usepackage[a4paper,left=2.2cm,right=2.2cm,top=2.1cm,bottom=2.2cm,
            headheight=12pt,footskip=1.0cm]{geometry}
\usepackage{microtype}
\usepackage{xltabular}
\usepackage{booktabs}
\usepackage{array}
\usepackage{ragged2e}
\usepackage{fancyvrb}
\usepackage{enumitem}
\usepackage{xcolor}
\usepackage{caption}
\usepackage{lastpage}
\usepackage{fancyhdr}
\usepackage{titlesec}
\usepackage[hidelinks]{hyperref}

% ---- body text ------------------------------------------------------------
\setlength{\parindent}{0pt}
\setlength{\parskip}{4pt plus 1pt}
\linespread{1.0}
\setlength{\emergencystretch}{2em}
\renewcommand{\arraystretch}{1.15}
\setlength{\tabcolsep}{5pt}
\captionsetup{font=small,labelfont=bf,skip=3pt}

% ---- headinds -------------------------------------------------------------
\titleformat{\section}
  {\normalfont\large\bfseries}{\thesection}{0.65em}{}
\titlespacing*{\section}{0pt}{13pt plus 2pt}{5pt}
\titleformat{\subsection}
  {\normalfont\normalsize\bfseries}{\thesubsection}{0.55em}{}
\titlespacing*{\subsection}{0pt}{10pt plus 1pt}{3pt}
\titleformat{\subsubsection}
  {\normalfont\normalsize\bfseries\itshape}{\thesubsubsection}{0.55em}{}
\titlespacing*{\subsubsection}{0pt}{8pt plus 1pt}{2pt}

% ---- running footer -------------------------------------------------------
\pagestyle{fancy}
\fancyhf{}
\fancyfoot[L]{\footnotesize FTEC5660 Reproducibility Study}
\fancyfoot[C]{\footnotesize Page \thepage\ of \pageref{LastPage}}
\fancyfoot[R]{\footnotesize Xiuling PIAO (1155267871)}
\renewcommand{\headrulewidth}{0pt}
\renewcommand{\footrulewidth}{0.4pt}
\fancypagestyle{plain}{%
  \fancyhf{}%
  \fancyfoot[L]{\footnotesize FTEC5660 Reproducibility Study}%
  \fancyfoot[C]{\footnotesize Page \thepage\ of \pageref{LastPage}}%
  \fancyfoot[R]{\footnotesize Xiuling PIAO (1155267871)}%
  \renewcommand{\headrulewidth}{0pt}%
  \renewcommand{\footrulewidth}{0.4pt}%
}

\hypersetup{
  pdftitle={Reproducing Mem0 on LOCOMO},
  pdfauthor={Xiuling PIAO (1155267871)},
  pdfsubject={FTEC5660 Reproducibility Work}
}

\begin{document}

\begin{center}
{\LARGE\bfseries Reproducing Mem0 on LOCOMO\par}
\vspace{0.35em}
{\large Xiuling PIAO \textperiodcentered{} 1155267871\par}
\vspace{0.25em}
{\small FTEC5660 Agentic AI for Business and FinTech \textperiodcentered{}
 Reproducibility Work\par}
\vspace{0.5em}
\rule{0.55\textwidth}{0.4pt}
\end{center}
\vspace{0.6em}

\textbf{Code: }\href{https://github.com/piupiu-star/mem0-repro}{\textbf{github.com/piupiu-star/mem0-repro}}\textbf{.}

All numbers in this report are generated by python -m src.summarize from the~per-question JSON files under results/; none is typed by hand.~results/summary.md is the machine-generated source of truth.

\section{Project summary and what I set out to reproduce}

\subsection{The System and the Claim}

\href{https://github.com/mem0ai/mem0}{Mem0} is a long-term memory layer for LLM agents: conversation turns are ingested into a memory store, relevant memories are retrieved per question, and an answerer conditions on them.~Ingestion is what makes this an agentic system. A single add call runs several LLM steps inside Mem0. It extracts candidate facts from the turn, compares them with memories already in the store, and then adds, updates or deletes entries. The caller supplies the conversation, and Mem0 decides what it keeps.

Mem0's README reports \textbf{92.5 \% LoCoMo accuracy} on its managed platform, with a headline of \textbf{>90 \% fewer tokens} and \textbf{91 \% lower p95 latency} than a full-context baseline (arXiv:2504.19413). LoCoMo (Maharana et al., arXiv:2402.17753) is a long-conversation QA benchmark: 10 conversations of 19--32 sessions each, questions split into multi-hop, temporal, open-domain, single-hop and adversarial.

\subsection{What I reproduced, modified and headline result}

The scope is one benchmark and one conversation, scored on all four categories:

\begin{center}\small
\begin{tabularx}{\textwidth}{@{}l>{\RaggedRight\arraybackslash}X@{}}
\toprule
\textbf{benchmark} & LOCOMO-10, snap-research/locomo, conversation 0 (19 sessions, reference date 22 Oct 2023) \\
\midrule
questions & first 120 in-scope questions: multi-hop 32, temporal 37, open-domain 13, single-hop 38 (adversarial excluded, as the authors do) \\
metric & LLM-as-a-Judge accuracy: the authors' verbatim judge prompt and parsing, but run on deepseek-chat rather than their default gpt-5 \\
reference numbers & README 92.5 \% overall + per-category; paper's relative token and latency claims \\
\bottomrule
\end{tabularx}
\end{center}

The modification is a single ablation that disables Mem0's LLM fact-extraction stage (add(..., infer=False)), so raw turns are stored verbatim instead of being rewritten into facts. It involves one flag, it is measurable through a paired test on identical questions, and it targets the stage the paper credits for its accuracy.

The headline result is 65.3 \% overall ($\pm$ 0.5 over three trials) against 92.5 \% reported. Two of the four categories beat the reported score, one is within 7 points, and one collapses. The collapse traces to a single reproducible cause in the open-source SDK, and once that category is set aside the rest of the benchmark matches. Excluding the temporal block, the remaining 83 questions score \textbf{229/249 = 92.0 \%} across three trials, against 92.5 \% reported. That is a measurement, not arithmetic: the 27.2-point headline gap is one category, one missing parameter, one SDK decision.

\section{Setup notes}

\subsection{env, data and keys}

Everything ran on one Windows 11 (25H2) laptop: Python 3.11.6, mem0ai 2.2.1, qdrant-client 1.20.0, fastembed 0.9.1 and spacy with en\_core\_web\_sm. No Docker, no Node and no GPU, which ruled the official evaluation scripts out. results/pip-freeze.txt records the freeze.

The corpus is locomo10.json, about 10 MB, downloaded at run time and kept out of the repository. Prompts and answer post-processing are copied verbatim from mem0ai/memory-benchmarks (Apache-2.0, attributed in src/prompts.py). The file also carries session\_N\_date\_time keys up to N = 35 with only 19 matching turn arrays, and src/locomo.py keys off the real arrays, so the reference date is 22 Oct 2023 rather than the orphan 4 Jan 2024.

The key sits in~.env, which git ignores. Only~.env.example~is tracked, and no key appears in any commit or result file.

\subsection{Compute and budget}

The harness records the token usage of every request it issues. At \$0.27 per million input and \$1.10 per million output tokens the study comes to about \$3.80, a floor rather than a total, because Mem0's extraction calls run inside~add()~and never reach the tracker. Ingesting the 419-turn conversation takes 535 s with extraction on and 10 s with it off.

\subsection{Model substitution}

The brief lets me swap the model provided the change is documented and held as a controlled variable. It is~deepseek-chat~through provider~deepseek, temperature 0.0 in every configuration, with the same token caps throughout. The embedder,~fastembed:BAAI/bge-small-en-v1.5, is a second departure forced on me because DeepSeek offers no embedding endpoint (§6, row 5). No prompt text was altered.

Results are not directly comparable to the paper because the underlying LLM differs. The embedder and the code version differ as well. A different answerer changes what is measured, a different judge changes the instrument, and the verbatim prompt fixes the rubric but not the grader, so every accuracy figure in section4 sets one model's reading against another's. Re-running the judge on gpt-5 is the first thing I would do with access to another model.

What \emph{is} held fixed is everything that defines the task and the rubric: the dataset, the question set, the retrieval depth, the judge prompt text, the answer post-processing and the arithmetic of the token counts.

\section{Reproduction targets and metric definitions}

\subsection{R1: absolute accuracy (primary)}

LoCoMo accuracy is the share of questions the judge labels CORRECT, under the authors' judge prompt and their rule that category-3 gold answers are cut at the first semicolon. The published figure is 92.5 \% overall, with 91.2 for single-hop, 91.3 for multi-hop, 72.7 for open-domain and 92.0 for temporal.~This number comes from the v3 README and was measured on Mem0's managed platform. The authors note that it includes proprietary optimizations unavailable in the open-source SDK, and that open-source users should expect ``directionally similar gains but not identical numbers''. R1 is therefore not a like-for-like target. It is still the number a reader reaches for, so the useful question is how wide the gap is and whether the direction survives.

I~did not run the authors' run.py. I wrote src/run.py instead, which repeats the same loop of ingest, search, answer and judge for each cutoff. Everything observable is held fixed: the judge system prompt and template are byte-identical copies, preprocess\_answer is copied, and the memory-rendering glue in the answer prompt was diffed against their get\_answer\_prompt, which surfaced one real mismatch that I then fixed (§6, row 9). What is not fixed is the file itself. This is a re-implementation and the diff is the only evidence that it matches, so the brief's requirement of the same evaluation script is met in substance rather than literally. I would rather say that than claim~a fidelity I never tested.

\subsection{R2 : relative efficiency (secondary).}

R2 covers the two efficiency claims, which need no judge and so cross fewer unknowns: over 90 \% fewer tokens and 91 \% lower p95 latency, both against full-context. Tokens per query means the answer prompt's prompt tokens, and the baseline puts all 419 turns into every prompt, as the paper's baseline does.

\subsection{R3 : version drift (a finding)}

The paper describes the v2 algorithm and repository main is v3, so an absolute number compared across those two lines measures a different system.

\subsection{Cost/latency instrumentation}

Search time is measured around memory.search() alone, separately from end-to-end time, which also includes answering and judging. The 0.88 s p50 and 91 \% p95 are search claims, and timing them another way would manufacture agreement that is not there.

\section{Results: my numbers against the reported numbers}

All figures below are at the authors' retrieval depth (top-200), pooled over three trials, and come from \texttt{results/summary.md}. Table 1 there also reports the top-10 curve and the spread across trials.

\subsection{Accuracy (R1)}

\begin{center}\small
\begin{tabular}{@{}llll@{}}
\toprule
\textbf{category} & \textbf{reported} & \textbf{measured} & \textbf{verdict} \\
\midrule
overall & 92.5 & 65.3 $\pm$ 0.5 & 27.2 points below \\
single-hop & 91.2 & 96.5 & 5.3 above \\
multi-hop & 91.3 & 84.4 & 6.9 below \\
open-domain & 72.7 & 97.4 & 24.7 above \\
temporal & 92.0 & 5.4 & 86.6 below \\
\bottomrule
\end{tabular}
\end{center}

Two categories beat the reported score, one is within 7 points, and one is destroyed. A single averaged number hides all of that, which is why I did not report one. Per-trial stability is $\pm$ 0.5 points overall, so the collapse is not a sampling accident.

\subsubsection{The temporal collapse, and its cause}

Temporal accuracy is 6 correct out of 111 question-instances, which is 2 out of 37 in each of the three trials, with no variation at all. That rules out noise, and it is not a limitation of the answering model either: the model answers confidently and wrongly because the memories it retrieves are wrong. One example, verbatim from~results/evidence/repro\_mem0\_ext\_temporal\_wrong.md: asked when Caroline went to the LGBTQ support group, the gold answer being 7 May 2023, the system retrieved a memory reading ``Caroline attended an LGBTQ support group on November 11, 2026, and found the experience very powerful'', generated that date, and was marked wrong for falling outside 14 days of the gold answer. The stored memory places the event three years in the future, in a conversation that took place in 2023.

The SDK source gives a chain of four steps.

The authors' LOCOMO runner anchors each session with~mem0.add(messages, user\_id, timestamp=session\_epoch)~(benchmarks/locomo/run.py, ingest loop). This is what allows relative expressions such as ``a few weeks before that'' or ``last Tuesday'' to be resolved against the conversation's own clock.

In~mem0ai==2.2.1, the open-source~Memory.add()~refuses that parameter:~timestamp~is gated in~mem0/memory/notices.py~and raises~ValueError: The timestamp parameter is not supported by the~OSS Memory SDK.~reference\_date~is gated identically.

My adapter catches the refusal and retries without~timestamp, so ingestion succeeds, but extraction now resolves every relative date against the wall clock of the machine, which at run time was November 2026.

Extraction then writes the wrong date into the memory text itself. The error is not recoverable at answering time: the prompt already forbids outputting 2025/2026, and the model obediently copies the date the store gave it.

That accounts for the whole headline gap. The paper's temporal result depends on a code path the open-source distribution does not permit. Nothing fails either: ingestion reports success, the store fills up, the run exits cleanly, and the only trace is one warning line the harness prints about the adapted call. Only the temporal questions show what it cost.

\subsubsection{Making the headline gap interpretable}

On the 83 questions the timestamp gate does not touch, I measure \textbf{229/249 = 92.0 \%} across three trials, against 92.5 \% reported. The 27.2-point headline gap is therefore not evidence of a broadly weaker system: it is one category, one missing parameter, one SDK decision.

\subsection{Efficiency (R2)}

\begin{center}\small
\begin{tabular}{@{}llll@{}}
\toprule
\textbf{claim} & \textbf{reported} & \textbf{measured (mem0, extraction on)} & \textbf{measured (full-context)} \\
\midrule
prompt tokens / query & 7.0 K & 10.1 K & 24.4 K \\
token saving vs full-context & >90 \% & 58.9 \% & --- \\
retrieval p50 & 0.88 s & 0.063 s & no retrieval step \\
retrieval p95 & --- & 0.141 s & no retrieval step \\
end-to-end p50 & --- & 3.55 s & 3.80 s \\
end-to-end p95 & --- & 5.12 s & 4.78 s \\
\bottomrule
\end{tabular}
\end{center}

\emph{Latency is quoted from trial 1 on both sides: the full-context baseline ran one trial while mem0 ran three, and comparing pooled mem0 latency against a single full-context trial would measure when each run happened, not the systems.}

\textbf{Token saving reproduces in direction but not in magnitude.} It is 58.9 \% against a claimed >90 \%, with mem0 still using 2.4$\times$ fewer prompt tokens than full-context.\textbf{The}\textbf{ 0.88 s p50 is not a like-for-like win.} My retrieval is 14$\times$ faster, but it runs against a local Qdrant file store on a laptop, not the managed platform.~\textbf{The 91 \% p95 reduction is not measurable as stated.} It is a reduction \emph{relative to full-context}, which performs no retrieval, so the ratio is zero-over-zero. Framed end-to-end, mem0 is \textbf{7.2 \% slower at p95} while being faster at p50 (3.55 s vs 3.80 s). Both configurations issue the same two calls, so this is tail behaviour, not call count: mem0's answer prompts vary with what was retrieved and the long ones cost more. A memory layer can be cheaper on average and still lose in the tail.

\subsection{Accuracy against full-context}

Full-context scores 99.2 \% (119/120, one trial, with no spread left to measure), against mem0's 65.3 \%. This rules the answering model out as the cause of the loss and states the trade plainly. Memory uses 2.4$\times$ fewer prompt tokens and gives up 34 points of accuracy, so the value proposition that survives here is cost, not accuracy

\subsection{Version drift (R3)}

The paper describes the v2 algorithm and the repository I ran is v3, while the 92.5 \% in the README is a v3 number from the managed platform. The comparison therefore crosses two gaps at once, v2 to v3 and platform to open source, and neither is separately observable from here. That is why I treat 92.5 \% as a reference point for the reader rather than a target to hit, and why the per-category pattern, with two categories above the reported level, one close to it and one destroyed, says more than the aggregate does.

\section{My modification and its measured effect}

\subsection{The change and Why this one.}

The change is~add(..., infer=False), which turns off Mem0's LLM fact-extraction stage so that each turn is stored verbatim rather than rewritten into facts. That stage is where the paper's accuracy comes from, since the memory layer decides what is worth keeping and compresses it into facts. §4.2 shows the temporal damage happens there as well, so turning it off is the smallest change that tests the mechanism directly. It is one boolean, set on every ingest call and on every fallback path, so an SDK signature change cannot quietly switch extraction back on.

\textbf{Result (top-200, three trials pooled, 360 question-instances).}~

Counts are shown as well as percentages, because at these sample sizes one question is worth 2.6 points in open-domain and 0.9 in single-hop:

\begin{center}\small
\begin{tabular}{@{}lllll@{}}
\toprule
\textbf{category} & \textbf{n} & \textbf{extraction ON} & \textbf{extraction OFF} & \textbf{change} \\
\midrule
overall & 360 & 235 = 65.3 \% & 279 = 77.5 \% & +12.2 \\
temporal & 111 & 6 = 5.4 \% & 41 = 36.9 \% & +31.5 \\
multi-hop & 96 & 81 = 84.4 \% & 89 = 92.7 \% & +8.3 \\
single-hop & 114 & 110 = 96.5 \% & 112 = 98.2 \% & +1.7 \\
open-domain & 39 & 38 = 97.4 \% & 37 = 94.9 \% & $-$2.5 \\
\bottomrule
\end{tabular}
\end{center}

The overall gain is \textbf{44 questions out of 360}. Thirty-five of them are temporal, eight multi-hop, two single-hop --- and open-domain moves by exactly \textbf{one} question in the wrong direction. At n = 39 that is noise, so I do not interpret it, and I do not present the modification as uniformly beneficial: it is a large gain in one category and a wash elsewhere.

\subsection{Paired test}

\begin{center}\small
\begin{tabular}{@{}lllllll@{}}
\toprule
\textbf{cutoff} & \textbf{baseline} & \textbf{modified} & \textbf{flips} & \textbf{delta} & \textbf{95 \% CI} & \textbf{McNemar p} \\
\midrule
top-200 (3 trials) & 235/360 & 279/360 & 51 gained / 7 lost & +12.2 pp & [+8.3, +16.4] & < 0.001 \\
top-10 (1 trial) & 74/120 & 76/120 & 8 gained / 6 lost & +1.7 pp & [$-$4.2, +7.5] & 0.791 \\
\bottomrule
\end{tabular}
\end{center}

The top-200 row pools three trials, so n = 360 is three repeats of the same 120 questions rather than 360 independent observations. Pooling inflates n and makes the p-value optimistic. The per-trial test is the conservative reading, and it agrees, giving +12.5, +12.5 and +11.7 pp. That is +12.2 $\pm$ 0.5, and the direction never reverses, so the effect is stable rather than a lucky trial.

The judge swap does not weaken this comparison. Both arms answer the same questions under the same DeepSeek judge, so the A/B delta compares two readings of one instrument. The judge confounder in §2 affects how far my score sits from the paper's 92.5 \%, not how far apart my two configurations are.

At the paper's operating depth the effect is established, and at top-10 it is not. That contrast is itself informative. The gain only appears when the answerer is allowed to see a broad slice of memory, which fits the date-in-the-text explanation, since a wrong date has to be retrieved before it can mislead.

\subsection{Why it helps, and how much of it I can attribute.}

Disabling extraction changes two things at once.

\textbf{Date anchoring.}~With extraction on, relative expressions are resolved against the wall clock and written into the memory text, so a 2023 conversation produces dates such as November 11, 2026. With extraction off, the original wording survives, and the answerer sees the conversation's own phrasing alongside its session dates.

\textbf{Granularity.}~503 extracted facts become 419 raw turns. Raw turns keep the connective phrases that extraction drops, such as before that or the same person, which multi-hop questions chain through. That is a plausible partial explanation for the eight extra multi-hop questions.

The per-category pattern points at channel 1, since temporal moves +31.5 while single-hop moves only +1.7. That is a hypothesis supported by the pattern, not an isolated causal estimate.~A clean separation would need a third arm that re-enables extraction and fixes the clock, and the OSS SDK does not let me construct one.

\subsection{Cost of the change.}

\texttt{infer=False} is not free:

\begin{center}\small
\begin{tabular}{@{}lll@{}}
\toprule
\textbf{} & \textbf{extraction ON} & \textbf{extraction OFF} \\
\midrule
ingest wall clock & 535 s & 10 s \\
memories stored & 503 & 419 \\
prompt tokens / query & 10.1 K & 11.4 K \\
token saving vs full-context & 58.9 \% & 53.3 \% \\
extraction LLM calls during ingest & yes --- inside add(), untracked & 0 \\
\bottomrule
\end{tabular}
\end{center}

The modification trades about 1.3 K tokens per query for 12.2 accuracy points and an ingest that runs 53 times faster. For a system whose selling point is cost, that is a trade-off rather than a free win, and it moves the token-saving claim further from the reported >90 \%.

\subsection{Limitation}

At top-200 the raw configuration retrieves 200 of its 419 stored turns, which is nearly half the conversation. Both configurations retrieve 200 items, but raw turns are longer, so the raw arm also puts about 13 \% more text into the prompt. Part of the gain could therefore come from seeing more text rather than from having cleaner memories. The top-10 column, where the gain disappears, is consistent with that reading, and I report it rather than choosing the depth that flatters the result.

\section{Debug diary}

Each row is something that broke on this machine, what it actually was, and what fixed it. Rows 6, 9 and 10 are the ones that changed my conclusions rather than just my code.

{\small
\begin{xltabular}{\textwidth}{@{}l>{\RaggedRight\arraybackslash}X>{\RaggedRight\arraybackslash}X>{\RaggedRight\arraybackslash}X@{}}
\toprule
\textbf{\#} & \textbf{Symptom} & \textbf{Root cause} & \textbf{Resolution} \\
\midrule
\endfirsthead
\toprule
\textbf{\#} & \textbf{Symptom} & \textbf{Root cause} & \textbf{Resolution} \\
\midrule
\endhead
\bottomrule
\endfoot
1 & RuntimeError: Storage folder \textasciitilde{}/.mem0/migrations\_qdrant is already accessed by another instance of Qdrant client & mem0 v3 opens a global telemetry Qdrant store, so a process may hold only one Memory instance & one configuration per process; enforced in src/run.py \\
2 & ValueError: Top-level entity parameters frozenset(\{'user\_id'\}) are not supported & mem0ai 2.x changed the search convention, and it raises ValueError, not TypeError --- a except TypeError fallback never fires & call search(query, filters=\{"user\_id": ...\}, top\_k=...), keep old forms as fallbacks \\
3 & Failed to load spaCy lemma model & the plain mem0ai wheel silently degrades v3's multi-signal retrieval to dense-only & install mem0ai[nlp] + en\_core\_web\_sm; now a documented requirement \\
4 & official OSS evaluation path is docker compose-only & the authors' harness wraps Mem0 in a FastAPI service & drive the SDK directly; deviations recorded in README.md §4 \\
5 & DeepSeek exposes no /embeddings endpoint & model substitution cannot be a single variable & local fastembed; recorded as a forced confounder \\
6 & temporal accuracy 2/20 while multi-hop scored 14/15 & timestamp= is refused by the OSS SDK (gated in mem0/memory/notices.py), so extraction anchored dates to the wall clock and wrote 2026 dates into memories & proven against the authors' runner, which does pass timestamp=session\_epoch; reported as the main finding rather than patched away \\
7 & ingest log printed stored=20 for a 419-turn conversation & Memory.get\_all() defaults to top\_k=20, so any bare call truncates and reports a fabricated total; the store behind that log held several hundred memories (503 in the run that produced the reported numbers --- §5) & pass top\_k=100\_000; cross-checked with QdrantClient.count(exact=True) \\
8 & run exited 1 although every artefact was written & PostHog/Qdrant shutdown noise on stderr, which PowerShell promotes to a terminating error & redirect to a log and judge success by artefact count, not exit code \\
9 & every memory was prefixed (2026-11-12) where the authors render (Thursday, November 12, 2026) & my date formatter had a shortened strptime tuple; mem0's ISO timestamps only match the authors' \%Y-\%m-\%dT\%H:\%M:\%S.\%f\%z candidate after [:26] truncation and \%z removal & copied \_to\_human\_date verbatim, verified identical output on 8 input shapes, archived every answer produced under the old prefix to results/\_archive\_v1\_datefmt/ and re-ran \\
10 & 120 questions scored against zero memories without raising & my own "never abort, score it 0" guard turned a hard failure into a plausible table & added a preflight that refuses to evaluate a store reporting 0 memories \\
11 & paper (v2) vs repository (v3) algorithms differ & the repo iterated after publication; the README score is a managed-platform v3 number & declared as R3; absolute numbers treated as a reference point, not a target \\
\end{xltabular}
}

Every number in §4 and §5 comes from the pipeline after all of the fixes above. Runs produced under the old behaviour are kept under~results/\_archive\_*/~and are not reported as results.

\section{Conclusions}

\subsection{What reproduces}

Single-hop reaches 96.5 \% against a reported 91.2 \%, open-domain reaches 97.4 \% against 72.7 \% on 13 questions, and multi-hop lands within 7 points at 84.4 against 91.3. The system is therefore not uniformly weaker than claimed, and it stays strong on exactly the question types that do not depend on dates. The direction of the efficiency claim also survives. Mem0 uses 2.4 times fewer prompt tokens than full-context, 10.1 K against 24.4 K per query, which is the substance of the claim that a memory layer beats stuffing the transcript. The pipeline itself holds up as well, since extraction, multi-signal retrieval, collection maintenance and answering all run end to end on a consumer laptop with no Docker, over a 419-turn conversation, for under two dollars.

\subsection{What does not reproduce}

\begin{center}\small
\begin{tabularx}{\textwidth}{@{}>{\RaggedRight\arraybackslash}Xl>{\RaggedRight\arraybackslash}X>{\RaggedRight\arraybackslash}X@{}}
\toprule
\textbf{claim} & \textbf{reported} & \textbf{measured} & \textbf{why} \\
\midrule
overall LoCoMo accuracy & 92.5 \% & 65.3 \% & dominated by the temporal collapse below \\
temporal accuracy & 92.0 \% & 5.4 \% & the OSS SDK refuses the timestamp anchor the authors' runner passes \\
non-temporal subset & 92.5 \% & 92.0 \% (229/249) & the 83 questions the gate does not touch \\
tokens per query & 7.0 K & 10.1 K & different embedder, different chunking, different platform \\
retrieval p50 latency & 0.88 s & 0.063 s & platform measurement vs local Qdrant file; not a like-for-like win \\
token saving vs full-context & >90 \% & 58.9 \% & direction holds, magnitude does not \\
p95 latency reduction & 91 \% & not measurable as stated; $-$7.2 \% end-to-end & full-context has no retrieval step, so the ratio is 0/0 \\
\bottomrule
\end{tabularx}
\end{center}

\subsection{Why}

Three causes, in order of how much they explain:

\noindent\hangindent=1.45em\hangafter=1 1. \textbf{A feature gate, not a bug.} Temporal accuracy does not degrade --- it is destroyed, because the anchor that makes dates resolvable is unavailable in the open-source package. The authors' runner passes \texttt{timestamp=session\_epoch}; \texttt{mem0ai==2.2.1} raises \texttt{ValueError: The timestamp parameter is not supported by the OSS Memory SDK.} Everything downstream inherits wrong dates, and the failure is silent at ingest. This is worth more than a matching percentage would have been: it tells a practitioner exactly which capability they lose by self-hosting.

\noindent\hangindent=1.45em\hangafter=1 2. \textbf{Platform against open source, across algorithm versions.}~The 92.5 \% is a managed-platform v3 number, the paper describes v2, and I ran open-source v3. That is three systems in one comparison. The authors' own README says as much when it warns that the scores include proprietary optimizations that open-source users do not get.

\noindent\hangindent=1.45em\hangafter=1 3. \textbf{Forced substitution.}~DeepSeek has no embedding endpoint, so the embedder changed too. Token counts and latency therefore cross more than one variable, which is why I lean on the per-category pattern and the paired test rather than on absolute agreement.

If the same harness is pointed at the \textbf{managed platform} with the authors' original models, the temporal block should recover and the headline should move towards 92.5 \%. If it does not, the platform-vs-OSS explanation is wrong and something else is doing the work. That is the single most informative follow-up, and it costs a platform subscription rather than more code.

\subsection{Recommendations to future users}

If your task involves dates, do not self-host Mem0 unpatched. Temporal accuracy is close to zero, and nothing warns you.

Never trust~Memory.get\_all()~for a count. It defaults to~top\_k=20, so a store of 500 is reported as 20. Use~QdrantClient.count(exact=True)~instead.

Dump per-question evidence, not just aggregate scores. Four of the eleven problems in §6 were silent failures that produced plausible numbers.

Check whether a latency claim is about retrieval or about the whole request.~A 91 \% reduction in retrieval latency is compatible with being slower end to end. Comparing the wrong pair produced a 98-point disagreement in my own table.

Measure full-context before adopting a memory layer. It scored 99.2 \% here at 2.4 times the tokens, so the trade depends on your budget and your accuracy floor, not on the memory layer's reputation.

infer=False~is a cheap, high-value ablation. It is a single flag, it makes no API calls, it speeds ingest up 53 times, and here it improved accuracy by 12.2 points (p~<~0.001).

\section*{Appendix: reproducing this report}

\begin{Verbatim}[fontsize=\footnotesize,frame=single,xleftmargin=4pt]
python -m venv .venv; .\.venv\Scripts\Activate.ps1
pip install -r requirements.txt
python -m spacy download en_core_web_sm
Copy-Item .env.example .env      # then add DEEPSEEK_API_KEY

python -m src.run --backend mem0 --infer 1 --conversations 0 --limit 120 --stage all
python -m src.run --backend mem0 --infer 0 --conversations 0 --limit 120 --stage all
python -m src.run --backend fullctx --conversations 0 --limit 120 --cutoffs 200 --stage all
.\tools\trials.ps1                # trials 2-3 at top-200
python -m src.summarize           # -> results/summary.md, results/runs.csv
python -m tools.evidence --tag repro_mem0_ext --category temporal --want wrong
\end{Verbatim}

Total: about 45 minutes of wall clock and roughly \$4 of API billing for the answer+judge stage, plus Mem0's untracked extraction calls during ingest.

\section*{References}

\noindent\hangindent=1.45em\hangafter=1 1. P. Chhikara, D. Khant, S. Aryan, T. Singh, D. Yadav. Mem0: Building Production-Ready AI Agents with Scalable Long-Term Memory. arXiv:2504.19413, 2025.

\noindent\hangindent=1.45em\hangafter=1 2. A. Maharana, D.-H. Lee, S. Tulyakov, M. Bansal, F. Barbieri, Y. Fang. Evaluating Very Long-Term Conversational Memory of LLM Agents. arXiv:2402.17753, 2024.

\noindent\hangindent=1.45em\hangafter=1 3. Mem0 repository and README, github.com/mem0ai/mem0 (accessed 13 Nov 2026). Source of the 92.5 \% platform figure, which the paper does not report.

\end{document}